\section{Translate the tree into a sequence}
\label{sec:tiny-stack-translation}

The constant folder changes a tree into another tree.
We can check its result with the same evaluator because the representation has not changed.
A compiler must eventually cross a harder boundary: the result and the input may have different shapes and different execution rules.
We need a way to explain why those different executions still agree.

Our first such translation removes the tree from execution.
It produces a sequence of instructions for a small stack machine.
The machine does not follow child pointers.
It reads instructions in order and keeps intermediate values on a stack.
The input still belongs to $V$, and every arithmetic result still belongs to $V$.
We change the execution mechanism without changing the arithmetic promise.

A \term{stack machine} obtains operands from a stack rather than naming a separate storage location for every intermediate value.
The top is the end from which values are removed.
For our notation, the top appears on the right.
Thus $[23,7,8]$ has $8$ at its top.
The sequence $[23,7]$ remains after removing that value.
This convention is small enough to state once and easy enough to get wrong if left unstated.

The instruction set has four operations.
\code{PUSH} carries a literal operand and pushes its value.
\code{GET\_INPUT} pushes the current input.
\code{ADD8} removes two values, adds them modulo $256$, and pushes the result.
\code{MUL8} does the same with multiplication.
There are no branch instructions and no explicit return instruction.
Reaching the end with exactly one stack value completes a program.
That value is its result.

The compiler's task is now concrete.
For every valid expression, it must produce instructions that finish with its value.
It must not consume a value that the expression's caller already had on the stack.
It must not rely on an instruction that reads a nonexistent operand.
These are related obligations, but they are not interchangeable.

\subsection{Give each instruction a meaning}
\label{sec:tiny-stack-semantics}

Write $S\mathbin{\|}[v]$ for the stack obtained by appending $v$ to stack $S$.
The stack $S$ can be empty.
For input $a$, the instruction rules are
\begin{align*}
  \code{PUSH}\ n &: S \longmapsto S\mathbin{\|}[n],\\
  \code{GET\_INPUT} &: S \longmapsto S\mathbin{\|}[a],\\
  \code{ADD8} &: S\mathbin{\|}[u,v]
    \longmapsto S\mathbin{\|}[(u+v)\bmod256],\\
  \code{MUL8} &: S\mathbin{\|}[u,v]
    \longmapsto S\mathbin{\|}[(u\times v)\bmod256].
\end{align*}
The arithmetic instructions need at least two values.
Trying to execute one with fewer values is \term{stack underflow}.
We reject such instruction sequences before running them.
This is an error in the target representation, not a newly introduced result of the source expression language.

For a binary instruction, $v$ is removed first because it is on top.
The value $u$ is removed second.
The calculation uses $u$ as its left operand and $v$ as its right operand.
Addition and multiplication give the same answer if those two operands are exchanged.
A later subtraction instruction would not.
We keep the operand convention explicit even though the present instruction set cannot expose that particular mistake through a changed numeric answer.

Let $\operatorname{Exec}(c,a,S)$ mean executing instruction sequence $c$ with input $a$ and initial stack $S$.
For now this notation describes the mathematical machine, with enough storage and a valid instruction sequence.
Execution of an empty sequence leaves $S$ unchanged.
Execution of a nonempty sequence applies its first instruction and then executes the remaining sequence.
The finite sequence and absence of jumps guarantee that this execution reaches its end.
The bounded C implementation will report failure when it cannot meet its storage contract.

\subsection{Generate instructions in operand order}
\label{sec:tiny-stack-generation}

Write $C(e)$ for the instruction sequence compiled from expression $e$.
The translation rules are
\begin{align*}
  C(\operatorname{lit}(n)) &= [\code{PUSH}\ n],\\
  C(x) &= [\code{GET\_INPUT}],\\
  C(\operatorname{add}(p,q)) &= C(p)\mathbin{\|}C(q)\mathbin{\|}[\code{ADD8}],\\
  C(\operatorname{mul}(p,q)) &= C(p)\mathbin{\|}C(q)\mathbin{\|}[\code{MUL8}].
\end{align*}
Here the same concatenation symbol joins instruction sequences rather than stacks.
The surrounding expressions identify which kind of sequence is involved.
These rules choose left-to-right operand evaluation.
That order is permitted because the source expressions are pure and total.
It is a choice made by this compiler, not a general permission to reorder effectful operations in other languages.

Algorithm~\ref{alg:tiny-stack-emit} gives the corresponding traversal.
Emission appends one instruction to a growable logical sequence.
The mathematical algorithm assumes enough storage; the implementation checks a fixed capacity.
Like folding, emission visits children before completing their parent.
Unlike folding, it does not replace the parent with another tree node.
It appends an instruction that will combine the values left by the children's code.

\begin{algorithm}
  \caption{Emit stack instructions for an expression.}
  \label{alg:tiny-stack-emit}
  \begin{pseudocode}
    \Require A valid expression tree $e$; an instruction sequence with enough capacity
    \Ensure Append $C(e)$ without changing the existing instruction prefix
    \Function{Emit}{$e$}
      \If{$e=\operatorname{lit}(n)$}
        \State append $\code{PUSH}\ n$
      \ElsIf{$e=x$}
        \State append $\code{GET\_INPUT}$
      \Else
        \State \Call{Emit}{$\operatorname{left}(e)$}
        \State \Call{Emit}{$\operatorname{right}(e)$}
        \State append $\code{ADD8}$ or $\code{MUL8}$ according to $e$'s operator
      \EndIf
    \EndFunction
  \end{pseudocode}
\end{algorithm}

Our running expression produces seven instructions.
The input and $7$ must be pushed before their addition.
The literals $3$ and $5$ must then be pushed before the second addition.
Finally, the multiplication combines the two subtree results.
Table~\ref{tab:tiny-stack-trace} works through the complete sequence at input $250$.
The first addition wraps to $1$, just as it did in the tree evaluator.

\begin{table}
  \centering
  \caption{The seven-instruction program produces 8 at input 250.}
  \label{tab:tiny-stack-trace}
  \begin{tabular}{rll}
    \toprule
    Position & Instruction & Stack after execution \\
    \midrule
    & Before execution & $[\,]$ \\
    0 & \code{GET\_INPUT} & $[250]$ \\
    1 & \code{PUSH} 7 & $[250,7]$ \\
    2 & \code{ADD8} & $[1]$ \\
    3 & \code{PUSH} 3 & $[1,3]$ \\
    4 & \code{PUSH} 5 & $[1,3,5]$ \\
    5 & \code{ADD8} & $[1,8]$ \\
    6 & \code{MUL8} & $[8]$ \\
    \bottomrule
  \end{tabular}
\end{table}

Figure~\ref{fig:tiny-stack-boundary} separates the representations involved.
The source evaluator and target executor reach the same value through different data structures.
The compiler does not run the target program to discover its answer.
It builds an instruction sequence that will work for every allowed input.
At input $9$, the same seven instructions produce $128$.
There is no input-specific substitution in the generated program.

\begin{figure}
  \centering
  \begin{tikzpicture}[node distance=15mm and 17mm]
    \node[basic block] (tree) {Expression tree};
    \node[basic block,right=of tree] (code) {Stack code};
    \node[basic block,below=of tree] (value1) {Value in $V$};
    \node[basic block,below=of code] (value2) {Value in $V$};
    \draw[cfg edge] (tree) -- node[above,font=\small] {compile} (code);
    \draw[cfg edge] (tree) -- node[left,font=\small] {evaluate} (value1);
    \draw[cfg edge] (code) -- node[right,font=\small] {execute} (value2);
    \draw (value1) -- node[above,font=\small] {equal} (value2);
  \end{tikzpicture}
  \caption{Correct translation connects two execution mechanisms at a common result.}
  \label{fig:tiny-stack-boundary}
\end{figure}

\section{Prove preservation of the surrounding stack}
\label{sec:tiny-stack-proof}
\index{stack machine!correctness}

Starting from an empty stack is how we run a whole expression program.
It is not enough for the most useful induction statement.
When the right subtree starts, the left subtree's value is already on the stack.
The right subtree's code must preserve that value while producing its own.
A proof restricted to an empty stack cannot directly justify that step.

Instead, we prove the stronger statement
\begin{equation}
  \operatorname{Exec}(C(e),a,S)
  = S\mathbin{\|}[\llbracket e\rrbracket_a]
  \label{eq:tiny-stack-prefix}
\end{equation}
for every valid expression $e$, every input $a\in V$, and every finite initial stack $S$ of values in $V$.
The execution must also avoid underflow.
Equation~\ref{eq:tiny-stack-prefix} says that compiled expression code preserves the entire initial stack and appends one result.
It states exactly the property that a parent needs from each child.

For a literal, the emitted instruction pushes that literal and changes nothing below it.
For the input node, the emitted instruction pushes $a$ and changes nothing below it.
These cases establish the statement for leaves.
Neither instruction removes an operand, so neither can underflow.

For an addition, assume the statement holds for children $p$ and $q$.
Executing $C(p)$ on $S$ gives $S\mathbin{\|}[u]$, where $u=\llbracket p\rrbracket_a$.
Now use the claim for $q$ with that larger initial stack.
Executing $C(q)$ gives $S\mathbin{\|}[u,v]$, where $v=\llbracket q\rrbracket_a$.
Both child executions avoid underflow by the induction assumptions.
The final \code{ADD8} has the two operands it needs.
It replaces them with $(u+v)\bmod256$ without touching $S$.
By the source addition rule, this value is $\llbracket\operatorname{add}(p,q)\rrbracket_a$.

For multiplication, the two child executions establish the same stack shape.
The final \code{MUL8} replaces $u,v$ with their product modulo $256$.
That is the source multiplication's value.
All expression forms have now been covered, so structural induction establishes Equation~\ref{eq:tiny-stack-prefix}.
Choosing $S$ empty gives the whole-program correctness statement.

The stronger statement was not extra decoration.
It let the proof follow the implementation's actual composition rule.
When an induction stalls because a recursive call runs in a larger context, the missing ingredient can be a claim about preserving that context.
Here the context is an existing stack prefix.
In a different compiler proof, it might be untouched memory, an environment of names, or a relation between source and target stores.
The appropriate claim depends on what the generated code can read or change.

\subsection{Do not confuse a fragment with a program}
\label{sec:tiny-stack-fragments}

Compiled expression code has a net stack effect of plus one.
That is true whether it starts with zero values or twenty.
A whole expression program additionally starts empty and finishes with exactly one value.
Those extra conditions belong to its execution interface.

For example, take a stack containing $23$ and execute the running expression's code at input $250$.
The result is $[23,8]$.
Nothing has gone wrong: the fragment preserved its caller's value and appended its result.
But if we present that two-value final stack as a complete expression program's result, we have violated the whole-program interface.

The C compiler can append a fragment to a code buffer that already contains instructions.
It therefore cannot promise that every successful append produces a complete executable program.
The existing prefix might leave values on the stack or might itself be malformed.
The compiler promises how it appends code; the verifier checks the resulting whole sequence.
Separating these promises prevents an append operation from being mistaken for an endorsement of arbitrary existing code.

\section{Check the code before executing it}
\label{sec:tiny-stack-verifier}
\index{stack machine!verification}

A target instruction sequence can have a plausible length and still be invalid.
An arithmetic instruction can appear before its operands.
A literal can lie outside $V$.
An instruction tag can be unknown.
The sequence can finish with zero values or leave several values instead of one.
We can detect all of these faults without knowing the runtime input.

The verifier scans the instructions while tracking only stack depth.
A push increases depth by one.
An arithmetic instruction requires depth at least two and decreases it by one.
At the end, depth must be exactly one.
The verifier also records the largest depth it encounters.
This is an exact bound for this branch-free machine: every accepted instruction executes on every run and its stack effect is independent of input values.

We also choose a canonical representation for unused instruction operands.
Only \code{PUSH} uses its argument field.
The argument field of every other opcode must contain zero.
This rule is not necessary for the arithmetic meaning, but it avoids several encodings of the same instruction caused by irrelevant leftover data.
It gives tests and diagnostic tools a simpler representation to inspect.
The compiler establishes the rule; the verifier checks it rather than trusting the compiler to have done so.

The verifier does not prove that the program implements a particular source tree.
The sequence consisting only of \code{PUSH} with operand $99$ is valid.
It is still the wrong translation of most source expressions.
Structural safety and semantic preservation remain different properties even when both are easy to test.

\subsection{Bound the storage explicitly}
\label{sec:tiny-stack-storage}

The implementation in \code{stack.h} reserves room for $64$ instructions.
This number is an engineering limit of the example, not part of the mathematical source language.
The code structure contains an instruction array and a logical length named \code{used}.
The valid portion is the prefix ending just before \code{used}.
The remaining slots are storage, not instructions in the program.

Because every push consumes an instruction, an accepted program of at most $64$ instructions cannot need more than $64$ stack slots.
The executor therefore uses a fixed stack array of that size.
The verifier rejects a logical code length larger than its instruction array before reading any instruction.
It checks every instruction before execution starts.
The executor then uses the verified stack effects to avoid an underflow or out-of-bounds stack write.

This argument assumes that the code object is a real, initialized C object and that its active instructions do not change between verification and execution.
It does not turn arbitrary bytes received from an untrusted file into a valid C structure.
There is no byte decoder in this example.
There is also no concurrent writer changing the instruction array during the run.
A serialized format or concurrent runtime would need additional boundary rules.

The compiler takes a limit no greater than the array capacity.
If an append would exceed the limit, it reports failure.
The public wrapper restores the original logical length when any recursive emission fails.
The existing active instruction prefix remains unchanged.
Unused slots may contain partially emitted instructions, but they lie outside the restored logical length.
This is a limited rollback guarantee, not a promise to restore every byte of the code object.

For the original seven-node example, an empty buffer with limit seven succeeds.
Limit six fails, and the buffer's logical length returns to zero.
After constant folding, the five-node expression needs five instruction slots.
Its code sequence pushes the input, pushes $7$, adds, pushes $8$, and multiplies.
The arithmetic result remains unchanged even though the capacity needed by this target representation falls.

\section{Read the compiler and executor together}
\label{sec:tiny-stack-c-implementation}

Listing~\ref{lst:tiny-stack-emitter} shows the recursive emission helper from \code{stack.c}.
The two recursive calls are separate statements.
Their left-to-right order is therefore visible in the implementation rather than depending on the evaluation order of function arguments in C.
The helper stops as soon as an append fails.
The public wrapper, shown in the source file, handles restoration of the logical length.

\begin{listing}
  \inputminted[firstline=13,lastline=25]{c}{example-code/ch01/stack.c}
  \caption{Emission follows the source tree while appending target instructions.}
  \label{lst:tiny-stack-emitter}
\end{listing}

Listing~\ref{lst:tiny-stack-runner} gives the executor.
It validates the input range and the whole code sequence before running any instruction.
The variable \code{sp} is a count of occupied stack slots.
Consequently, the next pushed value goes into \code{stack[sp]}, and removing a value first decrements \code{sp}.
Keeping that convention consistent is more useful than relying on a variable name such as \enquote{top}, which could mean either an index or a count.

\begin{listing}
  \inputminted[firstline=66,lastline=85]{c}{example-code/ch01/stack.c}
  \caption{The executor runs verified code and publishes one result on success.}
  \label{lst:tiny-stack-runner}
\end{listing}

The executor removes the right value into \code{b} and then the left value into \code{a}.
Its unsigned intermediate arithmetic has the same range argument as the tree evaluator.
Both operands lie in $V$, so the largest product is $65025$.
The explicit remainder implements the target instruction's specified arithmetic.
The implementation does not rely on a host overflow to stand in for the guest machine's reduction.

The output parameter is assigned only after successful verification and execution.
On rejection, the caller's existing result value remains unchanged.
Failure is reported through a Boolean status rather than a special numeric result.
Every value in $V$, including zero, is a legitimate program result.
Using zero to mean both a valid answer and a failed execution would make the interface needlessly ambiguous.

\subsection{Account for both time and space}
\label{sec:tiny-stack-costs}

For a tree with $N$ nodes and height $H$, successful compilation emits exactly $N$ instructions.
Each leaf emits one push and each binary node emits one arithmetic instruction.
With a preallocated instruction buffer and constant-time append, compilation takes $O(N)$ time and $O(H)$ additional recursion space.
The target code occupies $O(N)$ slots.
This statement describes the algorithm as capacity grows; the example's fixed limit restricts which instances it accepts.

The verifier takes $O(N)$ time and $O(1)$ auxiliary space when instruction count is $N$.
The executor takes another $O(N)$ time.
This implementation verifies on every call, so a reported execution cost includes both scans.
A runtime could verify an immutable code object once and execute it many times, but it would then need a trustworthy association between the checked object and the object being executed.
Skipping the scan without establishing that association would remove the premise of the stack-safety argument.

The exact peak stack demand can be calculated from the source tree.
Let $D(e)$ be the maximum number of extra stack slots needed while executing $C(e)$.
For either leaf, $D(e)=1$.
For either binary operation with children $p,q$,
\begin{equation}
  D(e)=\max\bigl(D(p),\,1+D(q)\bigr).
  \label{eq:tiny-stack-demand}
\end{equation}
The first term covers evaluation of the left child.
The second covers evaluation of the right child while the left result occupies one slot.
The final arithmetic instruction does not increase the peak.

For the original running expression, each addition has demand two.
The root therefore has demand $\max(2,1+2)=3$.
After folding the right addition into a literal, the demand becomes $\max(2,1+1)=2$.
The verifier obtains the same peaks from the instruction sequences without seeing the original trees.
Agreement between the two calculations is another useful implementation check.

\section{Compose folding with translation}
\label{sec:tiny-stack-composition}

We now have two distinct preservation results.
Folding preserves the source expression's value.
Compilation preserves that value while moving from tree evaluation to stack execution.
The two results can be connected because folding produces a valid tree of the kind compilation expects.
That representation condition is essential: equality of values would not excuse a malformed tree passed to the compiler.

With enough target capacity, the combined claim is
\begin{align*}
  \operatorname{Exec}(C(F(e)),a,[\,])
    &= [\llbracket F(e)\rrbracket_a]\\
    &= [\llbracket e\rrbracket_a].
\end{align*}
The first equality comes from the compiler's stack-prefix theorem.
The second comes from the folder's value-preservation theorem.
Neither proof needs to know the internal implementation of the other pass.
Their contracts identify the shared boundary.

The bounded implementation adds a status distinction to this account.
A translation that reports insufficient capacity has not produced an executable replacement.
The preservation claim concerns successful translations under the stated resource assumptions.
The caller must handle failure rather than execute unused or partially written array slots.
In a larger compiler, that handling might grow a buffer, select another representation, or report a controlled compilation failure.
It must not quietly reinterpret a resource failure as the source program's answer.

\subsection{Test the connected claim}
\label{sec:tiny-stack-testing}

The new test program, \code{test\_stack.c}, compares source evaluation with target execution before and after folding.
For each selected tree, it records all $256$ source results before mutation.
It compiles the original tree and compares the target result at every input.
It then folds the tree, compiles again, and compares the second target program against the saved source results.
Saving the original answers prevents the test from accidentally treating the already folded tree as its only reference.

The running example gets additional structural checks.
The test checks all seven original opcodes and their operands, not merely the final numeric result.
It checks peak depth three before folding and peak depth two afterward.
It checks that the folded program has five instructions.
Those assertions expose a translation that preserves answers but fails to implement the intended representation or simplification.

For the arithmetic instructions, the tests cover both operations on all $65{,}536$ literal operand pairs.
For translation across larger shapes, they generate $256$ bounded trees and check every input before and after folding.
Those generated cases remain a finite sample of source shapes.
Their value is in testing the connections among components, not in replacing the structural proof with a large count of successful runs.

Rejection tests exercise the target interface separately.
They include an empty program, an arithmetic instruction with no operands, an out-of-range literal, an unknown opcode, nonzero unused operands, excess final values, and an impossible logical length.
They also check that a failed run leaves its output parameter alone.
Capacity tests exercise partial emission and restoration of the logical prefix.
A nonempty-prefix test confirms that a successful fragment append need not constitute a valid whole program.

These tests cover a boundary that result comparisons alone would miss.
A compiler used only on valid small trees might never emit an invalid instruction.
Its verifier still needs rejection tests because the verifier's contract is broader than the compiler's normal output set.
Conversely, a verifier that accepts a program says nothing about whether the compiler chose the correct program for its source.
Testing the components together does not remove the need to test their separate promises.
